\chapter{稳健性检验补充表}
\label{chap:appb-robustness-tables}

\section{替代设定与补充结果}
\label{sec:appb-alternative-specifications}

本附录用于演示在中文册中追加第二个附录时，目录、书签、章节编号和表格编号是否仍能保持稳定。相关变量定义与构造方式已在附录~\ref{chap:appa-variable-notes} 中说明。表~\ref{tab:appb-robustness} 给出三个示例性的稳健性检验设定，展示监督与激励变量在不同控制组下的方向是否保持一致。

\begin{table}[htbp]
    \centering
    \caption{稳健性检验补充表示例}
    \label{tab:appb-robustness}
    \begin{tabular}{lccc}
        \toprule
        变量 & 设定一 & 设定二 & 设定三 \\
        \midrule
        monitor\_intensity & $-0.082^{**}$ & $-0.076^{**}$ & $-0.069^{*}$ \\
        incentive\_pay & $-0.115^{***}$ & $-0.109^{***}$ & $-0.101^{**}$ \\
        task\_complexity & $0.037$ & $0.042$ & $0.033$ \\
        企业固定效应 & 是 & 是 & 是 \\
        年份固定效应 & 是 & 是 & 是 \\
        \bottomrule
    \end{tabular}
\end{table}

从这个示例可以看到，即便控制变量和样本口径发生变化，监督强度与绩效激励的方向仍保持稳定。若正文需要引用补充回归，可直接指向附录表~\ref{tab:appb-robustness}。
